\scanpage{54}
% archive sheet number: 53
Sous ces conditions, il existe une partie
\margin{\rule[-5\baselineskip]{0.35pt}{5\baselineskip}}\deleted{fermée} $H$ de $G$ dont l'image inverse dans
$G'$ soit $H'$. (Donc si le morphisme
$S'\longrightarrow S$ \deleted{\unsure{fid.} de $S$ \unsure{un} \unsure{morph.} top. quotient}
\inserted{est top. submersif}, \deleted{\unsure{$H$}}
p. ex. est propre, ou plat avec $S$ loc. noeth)
alors $H$ est une partie fermée de $G$.
En tout cas, c'est ``ensemblistement un
ss-groupe''.

% slanted
\margin{\uline{Cor. 4} Sous les conditions du Cor. 3, supposons que $S=\unsure{\operatorname{Spec}|\operatorname{Spm}}(k)$, \unsure{avec} $k$ \uline{parfait}, \deleted{\unsure{$S'$ rigide}} \deleted{\illegible{1}} (\unsure{bien entendu|ceci entraîne} les conditions en question). Alors il existe un ss-groupe fermé $H$ de $G$, \unsure{simple sur $k$|réduit sur $k$}, tel que \deleted{\unsure{$H'=H\otimes_k S'$|$H'=H\times_S S'$}} \inserted{$H'\supset H\otimes_k S'$} \deleted{et que} \inserted{et que} $H'$ ait $=$ ensemble ss-jacent que $H\otimes_k S'$.}

Le corollaire 3 résulte facilement du
corollaire 2, en se ramenant au cas où
$S=\operatorname{Spec}(k)$, et où on prend pour $k'$ le corps
\deleted{\unsure{du pt}} résiduel \unsure{des|du} pt gén. de $S'$.

Démonstration de la proposition. \deleted{On peut
supposer que} Appliquant le th. de
rigidité à $H'_0$, on voit qu'il
existe un ss-groupe $H_0$ \deleted{\unsure{du|de}} $G$ tel que
$H'_0=H_0\otimes_k k'$. Passant au quotient par
$H_0$, on est ramené au cas où $H'$ est
un groupe \uuline{fini}. Il est donc annulé
par un entier $n$, donc contenu
dans $\,{}_nG'=({}_nG)\otimes_k k'$. Cela nous ramène
au cas où $G$ est fini sur $k$. \deleted{\unsure{La}}
